\documentclass[11pt,a4paper]{article}

\usepackage[utf8]{inputenc}
\usepackage[T1]{fontenc}
\usepackage{amsmath,amssymb,amsfonts}
\usepackage{physics}
\usepackage{geometry}
\usepackage{enumitem}
\usepackage{booktabs}
\usepackage{hyperref}

\geometry{margin=2.5cm}

\newcommand{\avg}[1]{\langle #1 \rangle}
\newcommand{\avgJ}[1]{\overline{#1}}
\newcommand{\re}{\operatorname{Re}}

\title{Justification of Gaussian i.i.d.\ couplings\\[4pt]
       in the $(2{+}4)$ random-laser model\\[2pt]
       \large Dense and sparse interaction graphs}
\author{}
\date{\today}

\begin{document}
\maketitle

%======================================================================
\section{Coupling integrals}
%======================================================================

In the stationary-regime Hamiltonian for a multimode random laser
\cite{Antenucci2015a,Gradenigo2020}, the four-body coupling is the
overlap integral of four normal-mode profiles weighted by the
third-order nonlinear susceptibility of the disordered medium:
\begin{equation}\label{eq:g4}
  g^{(4)}_{i_1 i_2 i_3 i_4}
  \;\propto\;
  \int_V \!d\vb{r}\;
  \sum_{\nu}
    E^{(\nu_1)}_{i_1}(\vb{r})\,
    E^{(\nu_2)}_{i_2}(\vb{r})\,
    E^{(\nu_3)}_{i_3}(\vb{r})\,
    E^{(\nu_4)}_{i_4}(\vb{r})\;
    \chi^{(3)}_{\nu}\!\bigl(\vb{r}\,\big|\,
      \omega_{i_1},\omega_{i_2},\omega_{i_3},\omega_{i_4}\bigr).
\end{equation}
The two-body coupling has the same structure with two mode profiles
and the linear susceptibility $\chi^{(1)}$:
\begin{equation}\label{eq:g2}
  g^{(2)}_{ij}
  \;\propto\;
  \int_V \!d\vb{r}\;
  \sum_\nu
    E^{(\nu)}_i(\vb{r})\,
    E^{(\nu)*}_j(\vb{r})\;
    \chi^{(1)}_\nu\!\bigl(\vb{r}\,\big|\,\omega_i,\omega_j\bigr).
\end{equation}
In a random medium, both the mode profiles $E_i(\vb{r})$ and the
susceptibilities $\chi^{(p)}(\vb{r})$ are spatially disordered
\cite{Fallert2009,Conti2008}.

%======================================================================
\section{Gaussianity of individual couplings}
%======================================================================

The integrals \eqref{eq:g4} and \eqref{eq:g2} can each be decomposed
into contributions from spatial cells $\{\Delta V_\ell\}$ of
linear size comparable to the spatial correlation length~$\xi$:
\begin{equation}\label{eq:cell_sum}
  g^{(p)}_{\bar\imath}
  = \sum_{\ell=1}^{M} \xi^{(p)}_\ell(\bar\imath),
  \qquad
  M \sim V/\xi^d \gg 1.
\end{equation}
Adjacent cells may be correlated, but cells separated by more than
$\xi$ are effectively independent.  Since the number of terms $M$ is
large, the Lindeberg--L\'evy central limit theorem (or, more
precisely, its extension to weakly dependent sequences
\cite{Billingsley1995}) guarantees that $g^{(p)}_{\bar\imath}$
converges in distribution to a Gaussian:
\begin{equation}
  g^{(p)}_{\bar\imath}
  \;\xrightarrow{\;M\to\infty\;}\;
  \mathcal{N}\!\bigl(\mu^{(p)},\;[\sigma^{(p)}]^2\bigr).
\end{equation}

\noindent
\textbf{Key point.}  This argument depends only on the ratio
$V/\xi^d$ and is entirely \emph{independent} of the topology of the
interaction network (fully connected, mode-locked, or sparse).
Whether a given quartet appears in the Hamiltonian is determined by
the frequency-matching condition (FMC) and, in the sparse variant,
by sub-sampling---neither of which depends on the spatial integral
that determines the coupling \emph{value}.  Hence Gaussianity of
each $g^{(p)}$ holds unchanged for dense and sparse graphs alike.

%======================================================================
\section{Statistical independence between couplings}
%======================================================================

Two couplings $g^{(4)}_{\bar\imath}$ and $g^{(4)}_{\bar\jmath}$
that share $k\ge1$ mode indices are in general correlated: their
overlap integrals involve the same global mode profiles
\cite{Zaitsev2010}.  The question is whether these correlations are
thermodynamically relevant.

\subsection{Dense graphs}

On the fully connected graph each mode participates in
$\binom{N-1}{3}=O(N^3)$ quartets.  The correlation between two
couplings sharing one mode decays as $O(1/V)$ from the spatial
integral; moreover, in the thermodynamic sum
$H_4=\sum g^{(4)}_{\bar\imath}(\cdots)$ the number of correlated
pairs is $O(N^3)$ per coupling out of a total of $O(N^4)$ terms.
Standard arguments from the theory of $p$-spin models
\cite{MPV1987,CrisantiSommers1995} then show that only the first
two moments of the coupling distribution enter the saddle-point
equations: correlations contribute corrections of order $O(1/N)$
to the free energy density.

After FMC filtering (mode-locked graph, $O(N^3)$ surviving quartets
\cite{Marruzzo2018}), every mode still participates in $O(N^2)$
quartets and the argument follows identically.

For $g^{(2)}$, $N(N-1)/2$ pairs with each mode in $O(N)$ of
them: correlations again contribute $O(1/N)$ corrections.

\subsection{Sparse graph: $O(N)$ sub-sampled quartets}

In the sparse variant, $n_s=N$ quartets are drawn uniformly at
random from the $O(N^3)$ FMC-surviving pool.  Each mode then
participates in $O(1)$ quartets on average (connectivity~$\sim 4$).
Two facts mitigate the loss of the $1/N$ suppression:

\begin{enumerate}[nosep,label=(\roman*)]
\item \textbf{Low collision probability.}
  The probability that two uniformly sampled quartets share at
  least one mode is
  $P(\text{shared}) = 1 - \binom{N-4}{4}\big/\binom{N}{4}
   \sim 16/N$,
  which vanishes for $N\to\infty$.  The overwhelming majority of
  coupling pairs in the sparse Hamiltonian involve disjoint
  index sets and are therefore exactly uncorrelated, regardless
  of the spatial statistics.

\item \textbf{Spatial decorrelation.}
  Even when two couplings do share an index, the correlation
  of the overlap integrals is $O(1/V)$ (the non-shared mode
  profiles wash out in the spatial average), which is negligible
  for macroscopic samples.

\item \textbf{Robustness of the RFOT scenario.}
  The random first-order transition (RFOT) characterising the
  glass phase of the $p$-spin model
  \cite{Kirkpatrick1989,Gradenigo2020} is known to be stable
  under changes in the coupling distribution---e.g.\ the random
  orthogonal model \cite{Marinari1994,Parisi1995} has strongly
  correlated couplings yet displays a qualitatively identical
  phase diagram.  This structural stability extends to sparse
  diluted models of the Viana--Bray type
  \cite{Viana1985,MPV1987}.
\end{enumerate}

For $g^{(2)}$ with uniform-frequency FMC, the number of surviving
pairs is $O(N)$ and each mode couples to $O(1)$ neighbours even
\emph{before} any sub-sampling.  The independence argument is thus
of the same nature as for the sparse $g^{(4)}$.

%======================================================================
\section{Summary}
%======================================================================

\begin{center}
\renewcommand{\arraystretch}{1.25}
\begin{tabular}{@{}lcc@{}}
\toprule
\textbf{Coupling / graph} & \textbf{Gaussianity} & \textbf{Independence} \\
\midrule
$g^{(4)}$, dense (FC or FMC) & CLT, exact for $V/\xi^d\gg1$
  & $O(1/N)$ corrections \\
$g^{(4)}$, sparse ($N$ quartets) & CLT, unchanged
  & collision prob.\ $O(1/N)$; spatial decorr.\ $O(1/V)$ \\
$g^{(2)}$, dense (FC) & CLT
  & $O(1/N)$ corrections \\
$g^{(2)}$, FMC ($O(N)$ pairs) & CLT
  & same regime as sparse $g^{(4)}$ \\
\bottomrule
\end{tabular}
\end{center}

\medskip\noindent
\textbf{Conclusion.}
The Gaussian approximation for each coupling is rigorously justified
by the central limit theorem applied to the spatial overlap integral,
and is valid for any interaction topology.
Statistical independence is exact in the thermodynamic limit for
dense graphs ($O(1/N)$ residual correlations).
On sparse graphs with $O(1)$ connectivity, the i.i.d.\ assumption
is not a mathematical identity, but is supported by three
independent mechanisms: (i)~the $O(1/N)$ probability that two
sub-sampled couplings share a mode index, (ii)~the $O(1/V)$ spatial
decorrelation of the overlap integrals, and (iii)~the proven
stability of the RFOT phase diagram under perturbations of the
coupling distribution.
These three mechanisms act concurrently and justify adopting
Gaussian i.i.d.\ couplings in all regimes of the $(2{+}4)$ model.

\bigskip
\begin{thebibliography}{99}

\bibitem{Gradenigo2020}
  G.~Gradenigo, F.~Antenucci, and L.~Leuzzi,
  Phys.\ Rev.\ Research \textbf{2}, 023399 (2020).

\bibitem{Antenucci2015a}
  F.~Antenucci, C.~Conti, A.~Crisanti, and L.~Leuzzi,
  Phys.\ Rev.\ Lett.\ \textbf{114}, 043901 (2015).

\bibitem{Antenucci2015b}
  F.~Antenucci, A.~Crisanti, and L.~Leuzzi,
  Phys.\ Rev.\ A \textbf{91}, 053816 (2015).

\bibitem{Angelani2006}
  L.~Angelani, C.~Conti, G.~Ruocco, and F.~Zamponi,
  Phys.\ Rev.\ Lett.\ \textbf{96}, 065702 (2006).

\bibitem{Leuzzi2009}
  L.~Leuzzi, C.~Conti, V.~Folli, L.~Angelani, and G.~Ruocco,
  Phys.\ Rev.\ Lett.\ \textbf{101}, 107203 (2009).

\bibitem{Fallert2009}
  J.~Fallert \textit{et~al.},
  Nat.\ Photonics \textbf{3}, 279 (2009).

\bibitem{Conti2008}
  C.~Conti and A.~Fratalocchi,
  Nat.\ Phys.\ \textbf{4}, 794 (2008).

\bibitem{Zaitsev2010}
  O.~Zaitsev and L.~Deych,
  J.~Opt.\ \textbf{12}, 024001 (2010).

\bibitem{Marruzzo2018}
  A.~Marruzzo, P.~Tyagi, F.~Antenucci, A.~Pagnani, and L.~Leuzzi,
  SciPost Phys.\ \textbf{5}, 002 (2018).

\bibitem{MPV1987}
  M.~M\'ezard, G.~Parisi, and M.~Virasoro,
  \textit{Spin Glass Theory and Beyond}
  (World Scientific, Singapore, 1987).

\bibitem{CrisantiSommers1995}
  A.~Crisanti and H.-J.~Sommers,
  J.\ Phys.\ I (France) \textbf{5}, 805 (1995).

\bibitem{Kirkpatrick1989}
  T.~R.~Kirkpatrick, D.~Thirumalai, and P.~G.~Wolynes,
  Phys.\ Rev.\ A \textbf{40}, 1045 (1989).

\bibitem{Marinari1994}
  E.~Marinari, G.~Parisi, and F.~Ritort,
  J.~Phys.\ A: Math.\ Gen.\ \textbf{27}, 7615 (1994).

\bibitem{Parisi1995}
  G.~Parisi and M.~Potters,
  J.~Phys.\ A: Math.\ Gen.\ \textbf{28}, 5267 (1995).

\bibitem{Viana1985}
  L.~Viana and A.~J.~Bray,
  J.~Phys.\ C \textbf{18}, 3037 (1985).

\bibitem{Billingsley1995}
  P.~Billingsley,
  \textit{Probability and Measure}, 3rd ed.\
  (Wiley, New York, 1995).

\end{thebibliography}

\end{document}
